\section{Ablation Study: Understanding the Role of Legal Context Components}
\label{sec:ablation_study}

To assess the individual contribution of each legal context component, factual narratives, statutory provisions, cited precedents, and semantically similar past cases, we perform an ablation study by systematically removing or altering these inputs across pipeline configurations. This study highlights how each component affects prediction accuracy and explanation quality, as reported in Tables~\ref{tab:prediction_pipeline_results} and~\ref{tab:explanation_pipeline_results}.

\paragraph{Impact on Judgment Prediction:}
The \texttt{CaseText + Statutes + Precedents} pipeline serves as the most comprehensive baseline. Removing statutory references (i.e., \texttt{CaseText + Precedents}) leads to a noticeable drop in F1-score (from 63.96 to 57.53 in the multi-label setting), indicating that legal provisions provide structured grounding essential for accurate predictions. Similarly, eliminating precedents (i.e., \texttt{CaseText + Statutes}) also reduces performance, though the drop is less steep, suggesting complementary roles of statutes and precedents. Pipelines relying solely on factual case narratives (e.g., \texttt{CaseFacts Only}) perform the worst, reaffirming that factual information alone is insufficient for robust legal outcome prediction.

\paragraph{Impact on Explanation Quality:}
A similar pattern emerges in explanation generation. The \texttt{CaseText + Statutes} pipeline consistently outperforms others across ROUGE, BLEU, METEOR, BERTScore, and G-Eval metrics, underscoring the importance of grounding explanations in explicit statutory language. When only precedents are added (without statutes), as in \texttt{CaseText + Precedents}, explanation scores drop significantly (e.g., G-Eval: 4.21 to 3.45 in the single-label case). The worst-performing setup is \texttt{Facts + Statutes + Precedents}, highlighting that factual inputs, even when supplemented with legal references, do not suffice for generating coherent and persuasive explanations if the core case context is missing.

\paragraph{Insights:}
These findings validate the design choices in \texttt{NyayaRAG}, where integrating factual case text with statutory and precedential knowledge mimics real-world judicial reasoning. Statutory references provide normative structure, while precedents offer context-specific analogies. Their absence not only reduces predictive performance but also degrades the factuality, clarity, and legal coherence of the generated explanations.

This ablation analysis also offers practical guidance: for retrieval-augmented systems deployed in legal contexts, careful curation and combination of retrieved statutes and relevant precedents are critical to ensure trustworthy outputs.
